\documentclass[11pt]{article}
\usepackage{amsmath}
\usepackage{amssymb}
\usepackage{geometry}
\geometry{margin=1in}

\title{Plan 01-01: J2 Potential in Spherical Coordinates}
\author{J2+J3+Drag Secular Perturbation Theory}
\date{2026-04-02 — Phase 1}

\begin{document}
\maketitle

\section{Objective}

Express the Earth's J2 gravitational perturbation potential in spherical coordinates and isolate the disturbing potential for perturbation analysis.

\section{Earth's Gravitational Potential with J2}

The gravitational potential of an oblate spheroid (Earth) including the dominant J2 zonal harmonic is:

\begin{equation}
U(r, \phi, \lambda) = -\frac{\mu}{r}\left[1 - J_2\left(\frac{R_E}{r}\right)^2 P_2(\sin \phi)\right]
\label{eq:full_potential}
\end{equation}

where:
\begin{itemize}
    \item $\mu = GM_\oplus = 3.986004418 \times 10^{14}$ m$^3$/s$^2$ — Earth's gravitational parameter
    \item $r$ — geocentric distance [m]
    \item $\phi$ — geocentric latitude [rad]
    \item $\lambda$ — longitude [rad] (does not appear due to axisymmetry)
    \item $R_E = 6.378137 \times 10^6$ m — Earth's mean equatorial radius
    \item $J_2 = 1.08263 \times 10^{-3}$ — second zonal harmonic coefficient [dimensionless]
    \item $P_2(x)$ — Legendre polynomial of degree 2
\end{itemize}

\section{Legendre Polynomial P2}

The second-degree Legendre polynomial is:
\begin{equation}
P_2(x) = \frac{1}{2}(3x^2 - 1)
\label{eq:P2}
\end{equation}

Therefore:
\begin{equation}
P_2(\sin \phi) = \frac{1}{2}(3\sin^2 \phi - 1)
\label{eq:P2_phi}
\end{equation}

\subsection{Special Values}

Verification at key latitudes:
\begin{align}
\phi = 0° \text{ (equator)}: \quad P_2(0) &= -\frac{1}{2}\\
\phi = 90° \text{ (pole)}: \quad P_2(1) &= 1\\
\phi = 45°: \quad P_2(\sin 45°) &= P_2(1/\sqrt{2}) = -\frac{1}{4}
\end{align}

\section{Disturbing Potential}

The \textbf{disturbing potential} (or perturbation function) $R$ is defined as the deviation from the point-mass potential:

\begin{equation}
R = U - U_{\text{point mass}} = U - \left(-\frac{\mu}{r}\right)
\end{equation}

Substituting Eq.~\eqref{eq:full_potential}:
\begin{equation}
R_{J2}(r, \phi) = -\frac{\mu}{r}\left[- J_2\left(\frac{R_E}{r}\right)^2 P_2(\sin \phi)\right]
\end{equation}

Simplifying:
\begin{equation}
\boxed{R_{J2}(r, \phi) = \frac{\mu J_2 R_E^2}{r^3} P_2(\sin \phi)}
\label{eq:disturbing_potential}
\end{equation}

Substituting $P_2$ from Eq.~\eqref{eq:P2_phi}:
\begin{equation}
R_{J2}(r, \phi) = \frac{\mu J_2 R_E^2}{2r^3}\left(3\sin^2 \phi - 1\right)
\label{eq:disturbing_expanded}
\end{equation}

\section{Dimensional Analysis}

Verify that $R_{J2}$ has dimensions of specific energy [m$^2$/s$^2$]:

\begin{align}
[R_{J2}] &= \frac{[\mu][R_E^2]}{[r^3]}\\
&= \frac{\text{m}^3/\text{s}^2 \cdot \text{m}^2}{\text{m}^3}\\
&= \frac{\text{m}^5/\text{s}^2}{\text{m}^3}\\
&= \text{m}^2/\text{s}^2 \quad \checkmark
\end{align}

This confirms $R_{J2}$ has dimensions of energy per unit mass, correct for a gravitational potential.

\section{Physical Interpretation}

\subsection{Sign of the Perturbation}

The disturbing potential $R_{J2}$ varies with latitude:

\begin{itemize}
    \item \textbf{Equator} ($\phi = 0°$): $P_2(0) = -1/2$ $\Rightarrow$ $R_{J2} < 0$ (attractive)
    \item \textbf{Poles} ($\phi = \pm 90°$): $P_2(\pm 1) = 1$ $\Rightarrow$ $R_{J2} > 0$ (repulsive)
    \item \textbf{Mid-latitudes} ($\phi = \pm 45°$): $P_2(\sin 45°) = -1/4$ $\Rightarrow$ $R_{J2} < 0$ (attractive)
\end{itemize}

This reflects Earth's oblate shape:
- The equatorial bulge creates additional gravitational attraction at low latitudes
- The polar flattening creates relative repulsion at high latitudes

\subsection{Axisymmetry}

The J2 potential depends only on $(r, \phi)$ and is independent of longitude $\lambda$. This axisymmetry is a consequence of Earth's approximate rotational symmetry about the polar axis.

\textbf{Implication:} The disturbing potential does not depend on the RAAN ($\Omega$) when expressed in orbital elements, which will lead to $\partial R_{J2}/\partial \Omega = 0$ and thus no secular change in inclination.

\section{Comparison with Point-Mass Potential}

The point-mass potential is:
\begin{equation}
U_{\text{point mass}} = -\frac{\mu}{r}
\end{equation}

The J2 correction is:
\begin{equation}
\frac{R_{J2}}{U_{\text{point mass}}} = -J_2\left(\frac{R_E}{r}\right)^2 P_2(\sin \phi)
\end{equation}

At LEO altitudes (h = 700 km, r ≈ 7078 km):
\begin{equation}
\left|\frac{R_{J2}}{U_{\text{point mass}}}\right| \approx J_2\left(\frac{6378}{7078}\right)^2 \approx 1.08 \times 10^{-3} \times 0.813 \approx 8.8 \times 10^{-4}
\end{equation}

The J2 perturbation is about 0.088\% of the point-mass potential — small but non-negligible for long-term orbit prediction.

\section{Success Criteria}

\begin{itemize}
    \item[$\checkmark$] J2 potential expressed with explicit Legendre polynomial: Eq.~\eqref{eq:full_potential}
    \item[$\checkmark$] Disturbing potential $R_{J2}$ isolated from point-mass term: Eq.~\eqref{eq:disturbing_potential}
    \item[$\checkmark$] Dimensional analysis confirmed: [m$^2$/s$^2$]
    \item[$\checkmark$] Physical interpretation provided
    \item[$\checkmark$] Axisymmetry noted
\end{itemize}

\section{Next Steps}

In Plan 01-02, we will express the geocentric latitude $\phi$ in terms of orbital elements $(i, \omega, f)$, transforming $R_{J2}$ from spherical coordinates to orbital element space.

\section{References}

\begin{itemize}
    \item Vallado, D. A. (2013). \textit{Fundamentals of Astrodynamics and Applications} (4th ed.). Equation 9-38.
    \item CONVENTIONS.md: Physical constants, Legendre polynomial definitions
\end{itemize}

\end{document}
